% ---------------------------------------------------------------------------
% Section 1: Introduction (input by main.tex; no preamble).
% Framing: the NEIS working paper on the Allied Industrial Base (neis2026allied; text in
% integration/aib_concept_review.md), paraphrased and cited, not copied.
% Numbers: final rerun of 2026-10-12, read from paper/results/*.json (method, es, eo, joint,
% coalition, strategic, bilateral, aib, aib_robustness, robustness).
% ---------------------------------------------------------------------------
\section{Introduction}\label{sec:intro}

Energy security used to be a matter of fuels: reserves, import routes and stocks of oil and gas.
As solar panels, batteries, grids and power electronics absorb a growing share of energy
investment, it is increasingly a matter of manufactured systems and of the minerals and components
they are made from \citep{carey2025secure, iea2024etp}. A country can now be secure in fuels and
exposed in transformers, battery materials or magnets. These manufacturing chains are
concentrated. In the trade data used in this paper, China is the largest world exporter in nine of
the 18 upstream and midstream technology--stage baskets that have a product definition, with
shares between 19 and 44 percent, and the refining of rare earths and the manufacture of magnets
and battery anode material are concentrated in China at shares close to or above 90 percent
\citep{iea2025gcmo, iea2026etp}. Concentration is not in itself a market failure; it reflects
scale economies, learning and geology. But it determines how far a disruption travels through
production networks \citep{acemoglu2012network, carvalho2021supply, baqaee2019macroeconomic},
and it turns trade into leverage. \citet{hirschman1945national} made the point for bilateral trade
dependence, and \citet{farrell2019weaponized} show that states with jurisdiction over the central
nodes of global networks can observe, restrict or deny access to them.

The response so far has been national. Governments have returned to industrial policy on a scale
not seen for decades \citep{juhasz2024new, evenett2024return}, often with sourcing rules that tie
subsidies to where minerals are processed and components made \citep{bown2023industrial}. National
programmes are meeting economic and political limits. Several economies have targeted the same
manufacturing segments without the demand or cost position to fill the plants, and announced
capacity has been cancelled as prices fell \citep{iea2026etp}; subsidy competition reaches allies
\citep{clausing2023carbon}; and tighter public finances make it implausible for every large
economy to fund a complete green-industrial agenda \citep{ladislaw2025evolution}. Smaller
economies cannot replicate the demand-led strategies of the largest ones at all
\citep{albertone2025regional}. Autarky is the expensive alternative. Simulations of broad
relocalisation find large losses of trade and output without consistent gains in resilience, and
diversification usually does more for resilience than reshoring \citep{oecd2025supply,
thakurweigold2024promoting}. Deep fragmentation into blocs is costly too, most of all for small
economies \citep{aiyar2023geoeconomic, javorcik2024economic}.

These observations underpin the proposition of an Allied Industrial Base
\citep{campbell2025underestimating, neis2026allied}. Matching China's scale across the
electro-industrial supply chains is beyond any one allied economy. The premise is that allied
economies jointly hold the markets, capital, resources and know-how that none holds alone, but
deploy them through competing national strategies \citep{campbell2025underestimating,
neis2026allied}. The working paper accordingly frames the policy goal as \emph{managed
interdependence}: each strategically important capability should be available at scale from some
trusted member, with redundancy where a failure would be very costly, specialisation where scale
economies and learning dominate, and open global sourcing where the risk is small
\citep{neis2026allied}. Implementing it requires answers to questions
that the literature mostly leaves open: which nodes of
which supply chains are severe enough to coordinate, who counts as a trusted partner, and which
member should produce what \citep[Section 5]{neis2026allied}.

This paper asks how a coalition of economies could allocate energy-technology supply chains
among its members, technology by technology and stage by stage, and builds the measurements that
such an allocation needs. The unit of analysis is a cell $\cell$: an economy, one of eleven
technologies (batteries, coal, electric grid, electric vehicles, gas, geothermal, green hydrogen,
nuclear, oil, solar and wind) and one of three stages (upstream extraction, midstream processing
and manufacturing, downstream deployment). Two pillars score the cells of 219 economies: Energy
Security ($\ES$), how secure an economy's supply is, and Economic Opportunity ($\EO$), how well
placed it is to compete as a producer. Both are built from rank-based S-curves of public data,
among them the IEA, the Energy Institute, the U.S. Geological Survey, UN Comtrade and the Atlas
of Economic Complexity (Section~\ref{sec:data}), and a cell
enters a ranking only if the economy is sovereign and observed data carry at least half of its
weight. A Strategic Index combines the pillars with a policy index into a within-country ranking
of sectors in which security enters as \emph{need}, favouring sectors where a country is exposed
but capable; an endowment screen keeps an extraction stage out of a member's priority unless the
member has the resource. The coalition is derived rather than chosen: from ideal points estimated
from UN General Assembly votes \citep{bailey2017estimating}, membership of U.S. trade and security
agreements, and shares of world supply. It has 31 core members on the U.S.-leaning side of the
Assembly with GDP of at least USD 100 billion and 26 partners admitted through agreements or as
major suppliers. Bilateral indices
score each partner as a friendshoring opportunity, where an economy's security need in a sector
meets the partner's economic opportunity on existing trade and investment ties, and as an export
market, taking account of alignment and, in four manufacturing stages, of the Chinese ownership of
the partner's foreign-invested plants. Finally,
an Allied Industrial Base analysis stress-tests 92 sub-sector nodes of the observed 2024 trade
network: it removes a dominant supplier, measures how much of the coalition's lost supply existing
relationships and optimized partnerships restore, simulates standardised capacity increments where
a gap remains, and groups the nodes that need joint action into Sectoral Compacts with exact
minimum coalitions. Figure~\ref{fig:framework_architecture}
sets out the architecture.

The paper makes four contributions. The first is measurement at the level of the supply-chain
node. Security and opportunity are scored jointly, on a common scale and with a common variable
architecture, for each stage of each technology; the scores decompose exactly into category and
variable contributions, and a rankability rule keeps cells that rest mostly on imputation out of
every ranking. Existing energy-security indices describe an energy system or a fuel, and
green-complexity measures describe products \citep{mealy2022economic, growthlab2024greenplexity};
neither pairs exposure with capability in the same cell. The second turns the pillars into
priorities while keeping the normative choices visible: the need reading of security, the
endowment screen, the basis of the policy index and the margin of each member's priority over its
runner-up are reported with the results. The third is a rule-based coalition and bilateral layer.
Alignment is measured from Assembly votes rather than asserted, the coalition and its tiers
follow from stated thresholds, and the bilateral indices carry both alignment and a first measure
of the common-mode risk that arises when suppliers in different countries share Chinese parents.
The fourth is a stress test of the allied supply network at the level of the sub-sector. It
separates the economies protected from those that may supply them, distinguishes what existing
relationships, coordinated trade and new capacity each contribute, checks trade-based capability
against physical production data where they exist, and forms Sectoral Compacts only where chains,
partners and policy needs cohere, with coalitions found by an exact binary program rather than by
ranking. It answers the working paper's question of where coordination pays \citep{neis2026allied}
with a measured coordination dividend, and it reports no composite score. The pipeline is open and configuration-driven, and the
corrections made during an end-to-end audit of an earlier version are documented in
Appendix~\ref{app:corrections}.

Five sets of results stand out. First, the coalition is stronger on security than on opportunity.
A coalition member leads 16 of the 30 security pairs and is in the top three of all 30, but the
best-placed member leads only 8 of the 33 opportunity pairs, against 22 for China. The two
pillars are correlated within pairs (median Spearman 0.59), largely through shared components:
on the 23 pairs in which both pillars keep components of their own, the median falls from 0.55
to 0.30 when only the non-shared components are used. About 15 percent of
the cells rankable on both pillars combine above-median opportunity with below-median security,
the region the Strategic Index favours. Second, members' priorities concentrate on manufacturing.
The index scores 3,154 cells in 174 countries, and every one of the 57 members has an allied
pick: 43 are midstream, and battery midstream alone is the pick of 27. The endowment screen
changes 41 picks, mostly away from upstream stages in which a member has a need but no resource.
The picks are fragile: the median margin over the runner-up is 0.015, 21 picks lie within 0.01,
and only five survive if security is read as strength rather than need. Third, alignment shapes
the partnership rankings. Viewed from the 31 core members, the United States and Germany are the
strongest partners and core members hold 24 of the top 25 places; giving the alignment component
zero weight would lift China from 44th to 3rd. Fourth, the allies' exposure is a capacity problem.
With China's supply removed, the 57 members lose 22.9 percent of their traded demand for
electro-industrial goods (47.5 percent in batteries); existing relationships and optimized trade
restore coverage only to 84.7 percent, and all of the recovery, USD 232 billion, is supply drawn
from buyers outside the coalition. New or expanded relationships add less than one point of coverage in
81 of the 91 nodes, and no node is classed as a case for trade coordination: 37 of the 70 electro-industrial nodes
need new capacity, and they form 22 Sectoral Compacts, all of them investment and offtake clubs, whose
exact minimum coalitions have three to nine countries. Fifth, rankings are more robust than
classifications. Country rankings within pairs barely move under random category weights (median
within-pair Spearman 0.997 for security and 0.988 for opportunity), and an empirical test of vertical
complementarities finds no out-of-sample gain from rewarding a stage for capability at adjacent
stages; the AIB classes, by contrast, depend on how the shock is read: if China's purchases fall with
its sales, 24 rather than 37 electro-industrial nodes need new capacity.

The results speak to some open questions of the working paper
\citep[Section 10]{neis2026allied} but not to others. It answers the question of when coordination
pays: against a China shock, rarely through trade alone, since the existing network already carries
nearly all that trade can do, and chiefly through capacity at capable hosts, organised by coalitions
of a few countries. It answers the question of how much duplication is enough only in part: a 25
percent supplier-share limit costs less than two points of coverage against the shock in 77 of the
91 nodes, but where capacity is missing no tested limit restores resilience to a second loss. The paper does not test the hypotheses that allied specialisation reaches a
resilience target at lower public cost than parallel national programmes, that demand coordination matters more than subsidy coordination, or that fencing
nodes dominates excluding sectors: these need cost, revenue and bankability data that the indices
do not contain (Section~\ref{sec:discussion}).

The rest of the paper is organised as follows. Section~\ref{sec:framework} sets out the
conceptual framework and notation and Section~\ref{sec:data} the data.
Section~\ref{sec:normalization} presents the normalisation, imputation and rankability rules.
Sections~\ref{sec:es} and~\ref{sec:eo} construct the Energy Security and Economic Opportunity
pillars, and Section~\ref{sec:joint} studies the joint space they define.
Section~\ref{sec:coalition} measures geopolitical alignment and derives the coalition,
Section~\ref{sec:strategic} builds the Policy and Strategic indices, Section~\ref{sec:bilateral}
the bilateral indices, and Section~\ref{sec:network} the Allied Industrial Base analysis of shocks,
partnerships and Sectoral Compacts.
Section~\ref{sec:robustness} reports robustness analyses, and Section~\ref{sec:discussion}
discusses the results, next steps and limitations.

% ---------------------------------------------------------------------------
\subsection{Related literature}\label{subsec:literature}

The Allied Industrial Base lies where several literatures meet that are usually read separately
\citep{neis2026allied}. We organise them by the streams that bear on its design and then turn to
the measurement literatures on which the indices build.

\paragraph{Industrial policy and capabilities.} \citet{rodrik2004industrial} recast industrial
policy as a process of discovery and strategic collaboration rather than the selection of
winners, and \citet{juhasz2024new} review a new empirical literature that is more favourable to it
while stressing context, measurement and governance; text-based and intervention-level data now
make policy measurable \citep{juhasz2022who, evenett2024return}. The green variant rests on
learning, coordination and environmental externalities \citep{rodrik2014green}, and its current
form ties support to location \citep{bown2023industrial, bistline2023economic}. This literature
studies the relation between a state and its firms. The coordination failures it identifies, such
as uncertain demand, missing suppliers and first-mover costs, do not stop at borders, and the
governance of industrial interdependence among several states remains thinly studied. Our Policy
index records the presence and density of policy by sector; it does not measure effectiveness.

\paragraph{Supply-chain resilience and economic security.} Production-network models show that
shocks to central suppliers need not average out \citep{acemoglu2012network,
carvalho2019production, baqaee2019macroeconomic}, that disruptions propagate along supply chains
\citep{carvalho2021supply}, and that private sourcing decisions can leave networks fragile
\citep{elliott2022supply}. \citet{grossman2023supply} analyse the choice between diversification
and reshoring, and \citet{baldwin2022risks} separate exposure from resilience. Policy work defines
a critical dependency as the combination of high disruption risk, high importance and limited
substitutability \citep{oecd2025economic}, and finds diversification generally more effective
than relocalisation \citep{oecd2025supply, thakurweigold2024promoting}. These results motivate
scoring by stage, separating exposure ($\ES$) from capability ($\EO$), and stress-testing the
allied network node by node. The AIB analysis shocks each node separately and does not model
propagation along the chain, and the bilateral indices capture
direct exposure only, not the indirect dependence that \citet{freund2024us} document.

\paragraph{Weaponized interdependence and allied scale.} \citet{farrell2019weaponized} show how
the hubs of global networks confer coercive power, which shifts attention from national market
shares to network position, a modern form of the concern of \citet{hirschman1945national}.
\citet{campbell2025underestimating} argue that only a coalition can match China's productive
scale and that alliances must pool capacity, not only political support. Friend-shoring
\citep{yellen2022remarks} and the evidence on fragmentation give this argument empirical content:
trade and investment between geopolitical blocs have weakened relative to flows within them
\citep{gopinath2025changing}, U.S. sourcing has moved away from China toward economies that are
themselves integrated with Chinese supply chains \citep{alfaro2023global, freund2024us}, and
confining trade to aligned partners carries welfare costs \citep{goes2022impact,
javorcik2024economic, aiyar2023geoeconomic}.

\paragraph{Regional industrial policy and historical coordination.}
\citet{albertone2025regional} show, for ASEAN, that cross-border industrial strategy must manage
competition as well as cooperation, and that coordination without integration needs mechanisms
that reveal synergies and monitor commitments. Historical cases point the same way. The Combined
Boards of 1942--45 pooled the raw materials and production of the United States, the United
Kingdom and Canada and allocated scarce output against total requirements
\citep{usdos1942combined, uswpb1944war}. Airbus combined national specialisation with shared
demand and public risk-taking \citep{ec2022industrial, imf2024industrial}. The European Coal and
Steel Community placed strategic sectors under common rules \citep{eu1951ecsc}. Each worked
around a concrete objective, clear roles and benefits reciprocal enough to sustain it.

\paragraph{Systemic bankability.} \citet{spokas2025systemic} show that infrastructure
technologies stall when any part of the surrounding system (offtake, permitting, enabling
infrastructure, insurance, supply chains, delivery capacity) is missing, so that a designated
capability is not a capability until it is bankable. Our indices do not measure bankability,
which is the main reason the investment step of the AIB analysis is a capacity simulation rather
than a project pipeline.

\paragraph{Competing schools.} The policy debate divides into three schools
\citep{neis2026allied}. The first defends multilateral openness and stresses the costs of
fragmentation \citep{aiyar2023geoeconomic, oecd2025supply}. The second, closest to the Allied
Industrial Base, accepts selective de-risking within managed interdependence. The third proposes
a reciprocal bloc that conditions market access on balanced obligations and the exclusion of
China \citep{cass2025grand}. Our coalition rule belongs with the second: it admits nonaligned and
even China-leaning economies through agreements or supply shares, and it fences off no sector as
a whole.

\paragraph{Measurement.} Four measurement literatures supply the building blocks. Energy-security
indicators organise availability, dependence, diversity and affordability
\citep{kruyt2009indicators, sovacool2011conceptualizing, ang2015energy, jewell2011iea} and define
security as the low vulnerability of vital systems \citep{cherp2014concept,
winzer2012conceptualizing}; mineral criticality combines supply risk with vulnerability
\citep{graedel2012methodology, nassar2020evaluating, kowalski2023raw}. Most such indices describe
an energy system or a fuel, whereas $\ES$ applies their logic to each stage of each technology.
Economic complexity treats exports as revealing capabilities \citep{balassa1965trade,
hidalgo2007product, hausmann2006structural, hidalgo2009building, hausmann2014atlas}, with
relatedness predicting future specialisation \citep{hausmann2022implied} and applications to
green products \citep{mealy2022economic, growthlab2024greenplexity}; $\EO$ combines these signals
with demand, cost, readiness and investment. The composite-indicator literature disciplines
normalisation, weighting and uncertainty analysis \citep{oecd2008handbook, saisana2005uncertainty,
paruolo2013ratings}, and our S-curve is the probability-weighting function of
\citet{karmarkar1978subjectively}. Alignment is measured with UN-vote ideal points
\citep{bailey2017estimating, voeten2009unga}, in the tradition of foreign-policy similarity
measures \citep{signorino1999taub} and with the bloc definitions of \citet{gopinath2025changing}.
The AIB layer draws on supply-chain network design and disruption-aware sourcing
\citep{geoffrion1974multicommodity, melo2009facility, tomlin2006value, snyder2016orms}, solved as
lexicographic linear programs \citep{ehrgott2005multicriteria}. It departs from that tradition by
restoring observed traded demand under explicit shocks rather than minimising cost, and by treating
the allied network as emergent from bilateral relationships rather than designed.

\paragraph{The gap.} Read together, these literatures diagnose the problem better than they
allocate a response \citep[Section 5]{neis2026allied}: they establish why capabilities and
concentrated dependencies matter, why allied scale can offset national limits and why autarky is
costly, but not who should produce what, or how much duplication suffices. Nor do they offer a
common measurement architecture that connects sector priorities, partner capabilities and
performance. The indices of \citet{carey2025secure} were a
first step. They score eleven technologies across commodities, manufacturing and deployment on
energy-security risk and economic opportunity, read the joint space as four strategic postures,
and rate partners as safer, stronger or more prosperous with weights 0.4, 0.2 and 0.4. This paper
takes them from rankings to a stress test of the allied network. It rebuilds the pillars at the level of the cell with
a rankability rule, derives the coalition from observed alignment and supply, rates every
directed partnership, and adds an analysis of observed trade that shows, node by node, where
existing relationships suffice, where coordination would help, where allies must build and which
countries a compact needs.
